\chapter{Number theory}

\section{GCD and Diophantine equations}
\begin{itemize}[nosep,leftmargin=*]
\item \texttt{gcd(a,b)}: nonnegative GCD; \texttt{gcd(0,0)=0}.
Absolute values must fit the common integer type.
Time $O(\log(\min(|a|,|b|)+1))$.
\item For $a,b\ge0$, LCM is zero if either is zero;
otherwise $(a/\gcd(a,b))b$.
\texttt{lcm(a,b)} requires a representable result.
Before multiplying, test $a/\gcd(a,b)\le\mathrm{MAX}/b$.
\end{itemize}
\kactlimport{ExtGCD.h}
\begin{itemize}[nosep,leftmargin=*]
\item $ax+by=c$, $(a,b)\ne(0,0)$, $g=\gcd(a,b)$:
solvable iff $g\mid c$. Scale B\'ezout coefficients by $c/g$.
All solutions from a particular $(x_0,y_0)$:
\[
 x=x_0+(b/g)t,\qquad y=y_0-(a/g)t,\quad t\in\mathbb Z.
\]
\item Bound integer $t$ with \texttt{floorDiv}/\texttt{ceilDiv}
(Mathematics), reversing inequalities for negative steps;
zero steps impose only feasibility. Widen intermediate products.
\item $a=b=0$: all pairs if $c=0$, otherwise no solution.
\end{itemize}

\section{Modular arithmetic}
\begin{itemize}[nosep,leftmargin=*]
\item m>0: normalize signed residues with \texttt{a\%=m; if(a<0) a+=m;}.
Residues do not preserve the order of the original values.
\item Product/power: \texttt{modMul}/\texttt{modPow} from Numerical's
\texttt{ModPower.h}. Cast before multiplying; ordinary / is not modular division.
\item Divide by b using its inverse, which exists iff $\gcd(b,m)=1$ for m>1.
For PRIME p and $p\nmid b$, inverse $b^{p-2}\bmod p$.
Reducing exponents modulo $\varphi(m)$ requires $\gcd(b,m)=1$.
\end{itemize}
\kactlimport{ModInv.h}
\kactlimport{GeneralCRT.h}
\kactlimport{ModLog.h}
\kactlimport{ModSqrt.h}
\kactlimport{ModSum.h}

\section{Primes and factorization}
\kactlimport{SPF.h}
\kactlimport{TrialDivision.h}
\kactlimport{Prime64.h}
\kactlimport{Factor64.h}

\section{Divisors}
\kactlimport{Divisors.h}
\begin{itemize}[nosep,leftmargin=*]
\item $n\ge1$, $n=\prod p_i^{e_i}$:
\[
 \tau(n)=\prod(e_i+1),\qquad
 \varphi(n)=n\prod(1-1/p_i).
\]
For $\varphi$, start at n and apply $v\mathrel{-}=v/p_i$
once per distinct prime.
\item Sum of divisors:
\[
 \sigma(n)=\prod(1+p_i+\cdots+p_i^{e_i})
 =\prod\frac{p_i^{e_i+1}-1}{p_i-1}.
\]
Use geometric sums modulo m unless denominators are invertible;
exact intermediate sums/powers may need \texttt{unsigned \_\_int128},
even for 64-bit n.
\item $n=1$: no prime factors; divisors $\{1\}$ and
$\tau(1)=\sigma(1)=\varphi(1)=1$.
\end{itemize}
\section{Totient and Mobius}
\kactlimport{TotientMobius.h}
\begin{itemize}[nosep,leftmargin=*]
\item $\sum_{d\mid n}\varphi(d)=n$, $\sum_{d\mid n}\mu(d)=[n=1]$.
For $\gcd(a,n)=1$, $a^{\varphi(n)}\equiv1\pmod n$.
\item Mobius inversion: if $g(n)=\sum_{d\mid n}f(d)$, then
$f(n)=\sum_{d\mid n}\mu(d)g(n/d)$.
For finite arrays, $g(n)=\sum_{n\mid d}f(d)$ gives
$f(n)=\sum_{n\mid d}\mu(d/n)g(d)$.
\item Primitive Pythagorean triples:
$(m^2-n^2,2mn,m^2+n^2)$ with $m>n>0$, coprime m,n of opposite parity.
Multiply by k>0 for all triples; interchange the legs if needed.
\end{itemize}
\par
